\documentclass{article}
\usepackage{hyperref}
\usepackage{Style}

\nocite{*} % Comentar si quiero citar
%\addbibresource{bibliografia.bib} % Quitar el comentado si quiero usar bibliografia

\begin{document}

\begin{minipage}{2.5cm}
    \includegraphics[width=2cm]{imagen_puc.jpg}
\end{minipage}
\begin{minipage}{12cm}
    {\sc Pontificia Universidad Católica de Chile\\
    Facultad de Matemáticas\\
    Departamento de Matemática\\
    Profesor: Manuel Cabezas -- Ayudante: Benjamín Mateluna}
\end{minipage}
\vspace{2ex}

{\centerline{\bf Introducción a la Geometría - MAT1304}
\centerline{\bf Ayudantía 11 - Repaso I4}}
\centerline{\bf 16 de abril de 2026}

\vspace{1cm}
\noindent\textbf{Problema 1.} Considere una circunferencia de centro $O$. Sean $\overline{AB}$ y 
$\overline{AC}$ tangentes al círculo en $B$ y $C$ respectivamente. Sea $\overline{CE}$ 
perpendicular al diametro $\overline{BD}$ con $E$ entre $O$ y $D$. Pruebe que 
$\overline{BE}\cdot\overline{BO}=\overline{AB}\cdot\overline{CE}$.

\vspace{5mm}
\noindent\textbf{Problema 2.} Considere dos circunferencias de centro $C$ y $D$, las cuales se 
intersectan en los puntos $A$ y $B$. Sea $P$ un punto sobre la circunferencia de centro $C$. Las 
rectas $\overleftrightarrow{PA}$ y $\overleftrightarrow{PB}$ intersectan en $R$ y $S$ 
respectivamente a la circunferecia con centro $D$. Muestre que el ángulo $\angle RPS$ es 
independiente de la elección de $P$.

\vspace{5mm}
\noindent\textbf{Problema 3.} Sea $\triangle ABC$ inscrito en una circunferencia de centro $O$.
Sea $D$ la intersección de la bisectriz de $\angle BAC$ con $\overleftrightarrow{BC}$ y $P$ la
intersección de $\overleftrightarrow{AB}$ con la perpendicular de a $\overleftrightarrow{OA}$ que
pasa por $D$. Demostrar que $\overline{AC}=\overline{AP}$.

\vspace{5mm}
\noindent\textbf{Problema 4.} Sea $P$ un punto dentro de una circunferencia, que no sea el centro. 
Dada una cuerda que pasa por $P$, que intersecta a la círculo en $A$ y $B$, muestre que
\begin{equation*}
    \overline{PA}\cdot\overline{PB}=r^{2}-d^{2}
\end{equation*}
Donde $r$ es el radio de la circunferencia y $d$ la distancia entre $P$ y el centro.

\vspace{1cm}
\noindent\textbf{\underline{Ejercicios Propuestos:}}

\vspace{5mm}
\noindent\textbf{Extra 1.} Sea $\triangle ABC$ inscrito en una circunferencia, considerar $M$
y $N$ en $\overline{AC}$ tales que $\angle ABM=\angle MBN=\angle NBC$. Sea $Q$ en la intersección
de $\overrightarrow{BM}$ y la circunferencia, por útimo, sea $P\in \overrightarrow{BN}\cap
\overline{QC}$. Demuestre que
\begin{equation*}
    \frac{\overline{AM}}{\overline{AN}}+\frac{\overline{CP}}{\overline{CQ}}=1
\end{equation*}
\textbf{Hint:} Vea que la expresión es equivalente a
\begin{equation*}
    \frac{PC}{PQ}=\frac{MN}{AM}
\end{equation*}

\vspace{5mm}
\noindent\textbf{Extra 2.} Sea $P$ un punto exterior a una circunferencia dada. Desde $P$ se
trazan dos rectas secantes a la circunferencia, que intersectan en los puntos $A,B$ y $C,D$. 
Muestre que
\begin{equation*}
    \overline{PA}\cdot\overline{PB}=\overline{PC}\cdot\overline{PD}
\end{equation*}

%\printbibliography % Quitar el comentado si quiero usar bibliografia

\end{document}